%% explanation_v3_non_changing.tex — single-ownship conflict-avoidance environment.
% !TEX program = pdflatex

\documentclass[11pt, a4paper]{article}
\usepackage{amsmath, amssymb}
\usepackage{booktabs}
\usepackage{geometry}
\usepackage{microtype}
\usepackage{parskip}
\usepackage{graphicx}
\usepackage[hidelinks]{hyperref}
\usepackage{algorithm}
\usepackage{algpseudocode}
\geometry{margin=2.5cm}

\title{\textbf{v3\_non\_changing Environment: Technical Specification}}
\author{}
\date{}

\begin{document}
\maketitle

% ─────────────────────────────────────────────────────────────────────────────
\section{Overview}
% ─────────────────────────────────────────────────────────────────────────────

The v3\_non\_changing environment is a deliberate simplification of
v3\_improved, designed to test whether a policy can learn basic separation
behaviour before tackling the full multi-aircraft control problem.

A single aircraft, referred to as the \textit{ownship}, is controlled by the
agent.
All other aircraft are uncontrolled \textit{traffic}: they receive no
instructions at any point and fly their initial headings throughout the
episode.
The reward is attributed solely to the ownship.
All sector, traffic density, and episode parameters are identical to
v3\_improved, so the two environments are directly comparable.

The sector area is drawn uniformly each episode from
10\,000--80\,000\,km$^2$.
Each episode contains $N \in [2,\,15]$ aircraft in total (ownship plus
$N-1$ traffic aircraft), where $N$ is determined by the sampled area divided
by a density drawn uniformly from 5\,000--15\,000\,km$^2$ per aircraft,
clipped to $[2,\,15]$.
All aircraft spawn at $M = 0.78$ (approximately 450\,kt TAS at FL350).

\paragraph{Key difference from v3\_improved.}
In v3\_improved the agent manages the whole sector: it selects the most
urgent aircraft each step and issues instructions to resolve conflicts between
any pair.
In v3\_non\_changing the agent manages only the ownship.
Traffic aircraft are passive objects; their trajectories are fixed by the
simulation physics and never modified.
This removes the focus-selection and commitment logic entirely, leaving a
clean single-aircraft MDP.

% ─────────────────────────────────────────────────────────────────────────────
\section{Environment Loop}
\label{sec:loop}
% ─────────────────────────────────────────────────────────────────────────────

\begin{algorithm}[h]
\caption{\textsc{Reset} — initialise a new episode}
\label{alg:reset}
\begin{algorithmic}[1]
\State Generate random convex polygon (5--7 vertices)
\State area $\sim\mathcal{U}(10\,000,\,80\,000)$\,km$^2$,\quad
       density $\sim\mathcal{U}(5\,000,\,15\,000)$\,km$^2$/aircraft
\State $N \leftarrow \operatorname{clip}\!\bigl(\lfloor\text{area}/\text{density}\rceil,\;2,\;15\bigr)$
\State Compute $T_\text{max}$ from sector diameter and cruise speed
\For{each slot $s = 1,\ldots,N$}
    \State Spawn aircraft on perimeter arc $s$ with uniform jitter
\EndFor
\State Ownship $\leftarrow$ aircraft in slot~0 \quad\textit{(fixed for the episode)}
\State \Return observation $\mathbf{o}(\text{ownship})$
\end{algorithmic}
\end{algorithm}

\begin{algorithm}[h]
\caption{\textsc{Step}(action) — advance one RL step (30\,s simulated)}
\label{alg:step}
\begin{algorithmic}[1]
\State Spawn any pending replacement aircraft
\If{action is heading offset $\delta$}
    \State $\psi_\text{cmd} \leftarrow \psi_\text{dest}(\text{ownship}) + \delta$
\ElsIf{action is Mach adjustment $\Delta M$}
    \State $M_\text{cmd} \leftarrow \operatorname{clip}(M_\text{cmd} + \Delta M,\;0.70,\;0.82)$
\EndIf
\State \textit{(Traffic aircraft receive no commands)}
\State Advance BlueSky simulation $60\times0.5$\,s $= 30$\,s
\For{each aircraft that exited the sector}
    \If{exiting aircraft is ownship}
        \State Apply wrong-direction exit penalty if $|\Delta\psi_\text{dest}| > 90^\circ$
    \EndIf
    \State Schedule replacement
\EndFor
\State Update ownship callsign if slot~0 was replaced
\State $r \leftarrow r_\text{safe} - w_\text{eff}\cdot\text{drift} - w_\text{work}\cdot c + r_\text{exit}$
\State \Return $\mathbf{o}(\text{ownship}),\; r,\;
               \text{truncated}=(\text{step\_count} \geq T_\text{max})$
\end{algorithmic}
\end{algorithm}

% ─────────────────────────────────────────────────────────────────────────────
\section{Action Space}
% ─────────────────────────────────────────────────────────────────────────────

Identical to v3\_improved.
Nine discrete actions: six heading offsets and three Mach adjustments, all
applied to the ownship only.
\begin{equation}
    \delta \in \{-30^\circ,\,-15^\circ,\,0^\circ,\,+15^\circ,\,+30^\circ,\,\varnothing\},
    \qquad
    \Delta M \in \{-0.04,\,-0.02,\,+0.02\}.
\end{equation}
Workload costs: 1.0 heading turns, 0.5 speed adjustments, 0 direct and hold.

% ─────────────────────────────────────────────────────────────────────────────
\section{Observation Space ($\mathbb{R}^{21}$)}
% ─────────────────────────────────────────────────────────────────────────────

The observation is constructed ego-centrically from the ownship and is
structurally identical to v3\_improved (Table~\ref{tab:obs}).
The four nearest traffic aircraft are included, sorted by predicted CPA
distance ascending.
Empty slots are padded with $(1,\,0,\,0,\,0)$.

\begin{table}[h]
\centering
\caption{Observation vector (21 features) — identical structure to v3\_improved.}
\label{tab:obs}
\smallskip
\begin{tabular}{clll}
\toprule
\textbf{Index} & \textbf{Feature} & \textbf{Definition} & \textbf{Range} \\
\midrule
\multicolumn{4}{l}{\textit{Ownship goal (2)}} \\[2pt]
0 & $\sin\Delta\psi$ & $\sin(\psi_\text{dest}-\psi_\text{cmd})$ & $[-1,1]$ \\[2pt]
1 & $\cos\Delta\psi$ & $\cos(\psi_\text{dest}-\psi_\text{cmd})$ & $[-1,1]$ \\[4pt]
\multicolumn{4}{l}{\textit{Ownship state (3)}} \\[2pt]
2 & Turn progress
  & $\operatorname{clip}((\psi_\text{cmd}-\psi_\text{actual})/30^\circ,\,-1,1)$
  & $[-1,1]$ \\[2pt]
3 & Time to exit
  & $\min(d_\text{boundary}/(v\cdot T_\text{look}),\,1)$
  & $[0,1]$ \\[2pt]
4 & Speed offset
  & $(M_\text{cmd}-M_\text{min})/(M_\text{max}-M_\text{min})$
  & $[0,1]$ \\[4pt]
\multicolumn{4}{l}{\textit{4 traffic aircraft $\times$ 4 features, sorted by CPA distance $\uparrow$:}} \\[2pt]
$+0$ & Distance  & $\min(d/3d_\text{sep},\,1)$                    & $[0,1]$  \\[2pt]
$+1$ & Forward   & $\cos(\beta-\psi_\text{own})$                   & $[-1,1]$ \\[2pt]
$+2$ & Right     & $\sin(\beta-\psi_\text{own})$                   & $[-1,1]$ \\[2pt]
$+3$ & CPA side  & $\operatorname{clip}(r^\text{right}_\text{CPA}/d_\text{sep},\,-1,1)$ & $[-1,1]$ \\
\bottomrule
\end{tabular}
\end{table}

% ─────────────────────────────────────────────────────────────────────────────
\section{Reward Function}
% ─────────────────────────────────────────────────────────────────────────────

The reward formula is identical to v3\_improved but evaluated for the ownship
only.
Traffic-vs-traffic conflicts carry no penalty.
\begin{equation}
    r_t = r_\text{safe}
        - w_\text{eff}\cdot\frac{1-\cos\Delta\psi_\text{own}}{2}
        - w_\text{work}\cdot c
        + r_\text{exit},
\end{equation}
\begin{equation}
    r_\text{safe}
        = -w_\text{safe}\cdot
          \max\!\left(0,\;1 - \frac{d_\text{CPA}^\text{min}}{2\,d_\text{sep}}\right),
\end{equation}
where $d_\text{CPA}^\text{min}$ is the minimum predicted separation between
the ownship and any traffic aircraft.
The exit penalty $r_\text{exit} = -w_\text{exit}$ applies only when the
ownship exits in the wrong direction.

\paragraph{Step-resolution limitation.}
With a 30\,s decision cycle, a fast head-on geometry (closure $\approx 900$\,kts)
can in principle produce a loss of separation that starts and fully resolves
within a single step.
In that case the reward is evaluated after the aircraft have already diverged,
so the worst-case separation is not directly penalised.
This is partially mitigated by the predictive $d_\text{CPA}^\text{min}$: once
past closest approach the code substitutes the current distance, which remains
small immediately after a breach.
The exponential reward shaping further reduces the impact by providing a
non-zero gradient in the steps preceding the conflict, so the policy is
penalised for converging geometry before any separation violation occurs.
The 30\,s step is retained because a shorter cycle reintroduces turn-aliasing,
making it impossible for the policy to observe the full geometric consequence
of a heading command.

\begin{table}[h]
\centering
\caption{Reward weights — identical to v3\_improved.}
\smallskip
\begin{tabular}{llr}
\toprule
$w_\text{safe}$  & Safety: CPA-margin penalty    & $1.00$ \\
$w_\text{exit}$  & Wrong-direction exit (sparse) & $0.50$ \\
$w_\text{eff}$   & Efficiency: heading drift     & $0.20$ \\
$w_\text{work}$  & Workload: instruction cost    & $0.05$ \\
\bottomrule
\end{tabular}
\end{table}

% ─────────────────────────────────────────────────────────────────────────────
\section{Comparison with v3\_improved}
% ─────────────────────────────────────────────────────────────────────────────

\begin{table}[h]
\centering
\caption{Key differences between the two environments.}
\smallskip
\begin{tabular}{lll}
\toprule
& \textbf{v3\_improved} & \textbf{v3\_non\_changing} \\
\midrule
Aircraft controlled & All $N$ (urgency-driven)  & Ownship only (slot~0) \\
Traffic instructions & Heading/speed as needed  & None \\
Focus selection      & Commitment rule          & Not applicable \\
Reward scope         & Focus aircraft per step  & Ownship always \\
Exit penalty         & Any aircraft             & Ownship only \\
Observation          & From most urgent aircraft & Always from ownship \\
Action space         & Discrete(9)              & Discrete(9) \\
Observation dim      & 21                       & 21 \\
\bottomrule
\end{tabular}
\end{table}

% ─────────────────────────────────────────────────────────────────────────────
\section{Training}
% ─────────────────────────────────────────────────────────────────────────────

Training uses the same PPO configuration as v3\_improved: four parallel
environments, network $[128,\,128]$, entropy coefficient $\beta_H = 0.02$,
and all other hyperparameters unchanged.
Results are saved to \texttt{Runs\_saved/v3\_non\_changing/}.

\end{document}
